\chapter{Fachlicher Hintergrund und Forschungsstand}
\label{ch:theorie}

\section{Begriffe und theoretisches Modell}
Ein idealisiertes Fadenpendel besteht aus einer punktförmigen Masse,
die an einem masselosen Faden befestigt ist. Für kleine Auslenkungen
lässt sich die Schwingungsdauer näherungsweise durch
\begin{equation}
  \label{eq:pendel}
  T = 2\pi\sqrt{\frac{l}{g}}
\end{equation}
beschreiben. Dabei ist \(l\) die Pendellänge und \(g\) die lokale
Fallbeschleunigung. Diese Kleinwinkelnäherung und ihre Voraussetzungen
werden etwa bei \textcite{openstax-pendel} erläutert.

\section{Gültigkeitsbereich und offene Fragen}
Das Modell setzt unter anderem eine kleine Amplitude und vernachlässigbare
Reibung voraus. Eine tatsächliche Untersuchung sollte festlegen, ob
diese Voraussetzungen im Versuchsaufbau hinreichend erfüllt sind.
% In der echten BLL hier: Quellenvergleich, Begriffsdefinitionen,
% dokumentierter Forschungsstand und Begründung des eigenen Beitrags.
